\subsection{Особенности и ограничения Bricks 2.3.0}\label{sec:bricks-limits}

Сгенерированный код каждого проекта проверяли. Вот что обнаружено.

\begin{center}\small
\begin{tabularx}{\linewidth}{@{}>{\raggedright\arraybackslash}p{0.3\linewidth}X@{}}
\toprule
\textbf{Ситуация} & \textbf{Что происходит и как обойти} \\
\midrule
Программа только из блоков светодиодов &
  блок светодиодов не добавляет в код \verb|pioneer = Pioneer()|; программа завершится с ошибкой \verb|name 'pioneer' is not defined|. \textbf{Обход:} поставьте в начало любой блок, читающий данные дрона (например, «напечатать» и «Получить данные от датчика расстояния») --- проекты 1, 2, 8 \\
«Распознать текст» вставлен прямо в «Получить текст из результата…» &
  генерируется \verb|:=| внутри генератора списка --- \verb|SyntaxError| в Python 3.12, программа не запускается. \textbf{Обход:} сохраните результат в переменную (проект 9) \\
Угол сервопривода &
  поле-циферблат округляет угол до шага $15^\circ$: 20 превращается в 15, $-40$ в $-45$ \\
Длинные блоки &
  вложенные блоки быстро становятся очень широкими; удобнее сохранять промежуточные результаты в переменные (проекты 7, 9, 10) \\
\bottomrule
\end{tabularx}
\end{center}

\begin{infobox}[Версия 2.3.0]
Список относится к Bricks~2.3.0 на Pioneer~OS 1.0.0; в других версиях поведение может отличаться. Если сомневаетесь, откройте сгенерированный \texttt{.py} и сравните с таблицей соответствия блоков.
\end{infobox}

\subsection{Блоки и Python: что с чем сравнивать}

Каждый из десяти проектов можно написать на Python из глав~\ref{sec:sdk2} и~\ref{sec:rknn}; это хороший способ закрепить и то, и другое.

\begin{center}\small
\begin{tabularx}{\linewidth}{@{}>{\raggedright\arraybackslash}p{0.38\linewidth}X@{}}
\toprule
\textbf{Проект Bricks} & \textbf{Где то же на Python} \\
\midrule
1, 2 --- светодиоды, циклы & светодиоды в главе~\ref{sec:sdk2} \\
3, 4 --- телеметрия, условия & дальномер и датчики в главе~\ref{sec:sdk2} \\
5 --- полёт по квадрату & полёт по точкам в главе~\ref{sec:sdk2} \\
6 --- сервопривод и фото & камера и сервопривод в главе~\ref{sec:sdk2} \\
7 --- найти людей & детекция \verb|Yolo| в главе~\ref{sec:rknn} \\
8 --- жесты & \verb|YoloPose| в главе~\ref{sec:rknn} \\
9 --- текст & \verb|PaddleOCR| в главе~\ref{sec:rknn} \\
10 --- ArUco & маркеры ArUco в главе~\ref{sec:python-vision}, с учётом кадра из \verb|Camera| \\
\bottomrule
\end{tabularx}
\end{center}

\begin{taskbox}[Из блоков в Python]
Выберите любой из проектов 1--4, откройте его сгенерированный \texttt{.py} и перепишите тело программы на Python вручную: без \verb|hex_to_bgr|, с обычными числами для цвета. Запустите в Pioneer Code.
\expected{ваша программа делает то же, что и блочная: те же цвета, паузы и вывод; \texttt{finally} с посадкой у вас на месте.}
\end{taskbox}

\begin{taskbox}[Свой проект]
Соберите в Bricks программу, которая показывает на светодиодах расстояние до препятствия: дальномер Ranger или датчик высоты, три цвета и пауза в полсекунды. Сохраните проект, экспортируйте его в \texttt{.xml} и импортируйте на соседний дрон.
\expected{цвет светодиодов меняется вместе с расстоянием; проект открывается на другом дроне без правок.}
\end{taskbox}

\begin{taskbox}[Где ошибка]
Соберите блок «Распознать текст» прямо внутри «Получить текст из результата…» и сохраните проект. Откройте сгенерированный код, найдите строку с \verb|:=| и исправьте проект через переменную \texttt{ocr}.
\expected{без переменной программа не запускается (\texttt{SyntaxError}), с переменной --- печатает распознанный текст.}
\end{taskbox}

\begin{missionbox}[Миссии симулятора к главе]
Режим \ui{Обучение} содержит трек «\ui{Мини 2 · Bricks}»: 16 миссий на блоках Pioneer Bricks. Миссия открывает дрон Пионер Мини~2 и редактор блоков Bricks, проверяет состав блоков и полёт в симуляторе. Миссии идут по проектам и логике главы: первые блоки и взлёт (1--2), светодиоды (проект~1, миссии 3--4), циклы и бегущий огонь (проект~2, миссии 5--6), данные и условия (проекты 3 и~4, миссии 7--9), полёт по точкам и квадрату (проект~5, миссии 10--12), камера, сервопривод и фото (проект~6, миссии 13--14), метки ArUco (проект~10, миссия 15) и итоговый световой патруль (16). Проекты 7--9 с нейросетями (люди, жесты, текст) и блоки с пометкой «только на дроне» в трек не вошли.
\end{missionbox}
